% ECONOMICS_PROBLEM: {"id": "F31-327", "title": "Uncovered-interest-parity deviations", "jel_code": "F31", "source_id": "OP-0332"}
% !TeX program = xelatex
\documentclass[11pt,a4paper]{article}
\usepackage[margin=23mm]{geometry}
\usepackage{fontspec}
\setmainfont{Latin Modern Roman}
\usepackage{amsmath,amssymb,mathtools}
\usepackage{xcolor,enumitem,booktabs,longtable,array,xurl,hyperref}
\definecolor{muted}{HTML}{53636C}
\hypersetup{colorlinks=true,urlcolor=blue,linkcolor=blue}
\setlength{\parindent}{0pt}
\setlength{\parskip}{5pt}
\newcommand{\R}{\mathbb R}
\newcommand{\E}{\mathbb E}
\newcommand{\N}{\mathbb N}
\newcommand{\ind}{\mathbf 1}
\newcommand{\Var}{\operatorname{Var}}
\newcommand{\Cov}{\operatorname{Cov}}
\newcommand{\supp}{\operatorname{supp}}
\DeclareMathOperator*{\argmin}{arg\,min}
\DeclareMathOperator*{\argmax}{arg\,max}
\newcommand{\entrymeta}[3]{{\sffamily\small\color{muted}\textbf{\texttt{#1}}\quad|\quad #2\par #3\par}}
\newcommand{\Problemref}[1]{\texttt{#1}}
\newcommand{\JELref}[2]{\texttt{#2} (JEL \texttt{#1})}
\begin{document}
\section*{Uncovered-interest-parity deviations}
\textbf{Problem ID:} \texttt{F31-327}\quad \textbf{JEL:} \texttt{F31}\par
% BEGIN ECONOMICS STATEMENT
\label{problem:OP-0332}
{\sffamily\small\color{muted}Original subject group: International finance and exchange rates\par}
\label{definitions:problem:30007614}
\textbf{Audit classification:} Background; proposed extension. Fama supplies conditional evidence about forward premiums; the recent IMF source supplies FX identification/policy background. The exact risk/belief/constraint decomposition is not verified as a source-stated open problem. Log and gross parity are distinguished.
\subsection*{Definitions and precise research target}
Let \(s_t\) be log domestic-currency units per foreign currency, and \(i_t,i_t^*\) comparable one-period domestic and foreign log gross interest rates. Define the log return difference from borrowing domestically and investing abroad as \(x_{t+1}=i_t^*-i_t+s_{t+1}-s_t\). Let \(\mathbb P\) be the objective probability law and \(\mathcal F_t\) dated available information. Under this conditional law, its expectation is the currency premium \(\pi_t=E_{\mathbb P}[x_{t+1}\mid\mathcal F_t]\); the log-parity convention requires \(\pi_t=0\); exact gross-return parity is separately defined below. An investor's subjective law \(\mathbb Q_t\) gives an expectation wedge
\[b_t=E_{\mathbb Q_t}[s_{t+1}\mid\mathcal F_t]-E_{\mathbb P}[s_{t+1}\mid\mathcal F_t].\]
The task is to identify which part of predictable excess returns is attributable to compensation for priced risks, belief errors, or trading and financing constraints. Specify a discount-factor model, a survey measurement model for \(\mathbb Q_t\), and the mechanism relating intermediary constraints to attainable portfolios. Determine the identifying restrictions separating these components and derive bounds when objective expectations or marginal investors are unobserved. Evaluate variation by currency, horizon and financial state. Realized excess returns may remain predictable under rational expectations with time-varying risk compensation; a regression slope inconsistent with parity does not by itself identify irrationality. Exact gross-return parity and log parity are distinct; use the chosen convention consistently.

\textbf{Additional formal definitions and scope (editorial).} Specify gross domestic and foreign rates \(R_t^d,R_t^f>0\), exchange rate \(S_t=e^{s_t}>0\), and net currency payoff \(X_{t+1}=R_t^f S_{t+1}/S_t-R_t^d\) per unit initially invested. Exact expected-return parity is \(E_t[X_{t+1}]=0\). If \(i_t=\log R_t^d\) and \(i_t^*=\log R_t^f\), then \(x_{t+1}=i_t^*-i_t+\Delta s_{t+1}\) is the exact log return difference; \(E_t[x]=0\) is a log-parity convention and differs from exact gross parity by Jensen terms. Under unconstrained common-belief pricing, \(E_t[M X]=0\), so \(E_t[X]=-\operatorname{Cov}_t(M,X)/E_t[M]\). With tastes, constraints or different subjective laws the Euler equation changes. Define the survey kernel and marginal investor before interpreting \(b_t\) as priced belief error.

For the identification part, declare an admissible class \(\Theta\) of structural models, including preferences, technologies, shock laws, measurement/sampling rules and any equilibrium selector. If \(O\) denotes the complete observable history and \(T(\theta)\) the target defined above, the identified set is \(\mathcal I_T=\{T(\theta):\theta\in\Theta,\ \mathcal L_\theta(O)=\mathcal L_{\rm obs}(O)\}\); point identification means this set is a singleton. A \(do(a)\) comparison replaces stated policy/technology equations while fixing the declared exogenous law and re-solving the remaining equations. These definitions do not assert that an identification theorem holds without further restrictions.
\subsection*{Variants and assumption changes}
\begin{enumerate}
\item \textbf{Rational expectations and exact gross parity (editorial).} Fix a common pricing kernel and identify risk compensation for \(X\), retaining Jensen effects when reporting logs.
\item \textbf{Surveys and financing constraints (editorial).} Add a measurement model for subjective expectations and a financing wedge; identify or bound their separate contributions rather than inferring beliefs from a forward-regression slope.
\end{enumerate}
\subsection*{References and source audit}
\begin{enumerate}
\item Primary journal scan verifies Eugene F. Fama, 1984, Journal of Monetary Economics 14:319--338; DOI matches journal PII. Locator: Introduction, p.319; Section 2, pp.320--323. Support: Forward-premium evidence under maintained assumptions. Full PDF accessible. URL: \url{https://eml.berkeley.edu/~craine/EconH195/Fall_14/webpage/Fama_Forward%20Discount.pdf}.
\item Official IMF/elibrary indexed records verify five authors, 2026 SDN/2026/003 and DOI. Official release is September 17, 2026; URL contains September 16. Locator: Official indexed abstract and currency-premium/decomposition passages; no exact full-page locator certified. Support: Fundamental/financial FX identification and policy background. Landing page initially opened; later full-text fetches failed. Indexed primary passages inspected only. URL: \url{https://www.imf.org/en/publications/staff-discussion-notes/issues/2026/09/16/drivers-of-exchange-rates-in-emdes-implications-for-foreign-exchange-intervention-578715}.
\end{enumerate}
\begin{enumerate}\raggedright
\item\hypertarget{definitions:source:30007614:1}{} Eugene F. Fama. 1984. \emph{Forward and Spot Exchange Rates}. Introduction, p.319; Section 2, pp.320--323.
\url{https://eml.berkeley.edu/~craine/EconH195/Fall_14/webpage/Fama_Forward%20Discount.pdf}.
DOI: \url{https://doi.org/10.1016/0304-3932(84)90046-1}.
\item\hypertarget{definitions:source:30007614:2}{} Ece Ozge Emeksiz; Andres Fernandez; Nikhil Patel; Ivan Petrella; Tatjana Schulze. 2026. \emph{Drivers of Exchange Rates in EMDEs: Implications for Foreign Exchange Intervention}. Official indexed abstract and currency-premium/decomposition passages; no exact full-page locator certified.
\url{https://www.imf.org/en/publications/staff-discussion-notes/issues/2026/09/16/drivers-of-exchange-rates-in-emdes-implications-for-foreign-exchange-intervention-578715}.
DOI: \url{https://doi.org/10.5089/9798229055208.006}.
\end{enumerate}

\subsection*{Common conventions: Source mathematical conventions}

These conventions apply to each entry unless explicitly replaced.

\textbf{Models and identification.} A primitive model \(m\) specifies all probability laws, feasible choices, information, equilibrium and measurement rules used by its target. The admissible class \(\mathcal M\) is fixed independently of outcomes. If \(P_O(m)\) is its observable probability law and \(\tau(m)\) the target, its identified set at observable law \(P\) is
\[
 \mathcal I_\tau(P)=\{\tau(m):m\in\mathcal M,\ P_O(m)=P\}.
\]
Point identification means that this set is a singleton. An empty set signals incompatibility of data and assumptions. Coordinate bounds need not describe the joint set. Bounds are sharp only if every retained target is attainable or arbitrarily approachable by compatible models. Finite bounds require explicit restrictions on unobserved quantities; unrestricted confounding, hidden wealth, extrapolation or arbitrary equilibrium selection can yield uninformative sets.

\textbf{Causal interventions.} A potential outcome \(Y_i(p)\) is the outcome under a complete policy regime \(p\), including timing, financing, information and market spillovers. The target population and units are fixed before treatment. With interference, use the complete assignment vector or a justified exposure mapping; individual no-interference notation alone cannot identify an economy-wide intervention. A difference of mean potential outcomes does not require the cross-world joint distribution. Conditioning on post-treatment migration, survival, enrollment or employment changes the estimand and needs separate justification.

\textbf{Statistical statements.} An estimator, confidence set, forecast or test specifies a sampling law, model class and loss/coverage criterion. Uniform \((1-\alpha)\)-coverage means \(\inf_{m\in\mathcal M}\Pr_m\{\tau(m)\in C_n\}\geq1-\alpha\), or an explicitly asymptotic version. The whole parameter space trivially has coverage, so informativeness requires a stated width, risk or power objective. A forecasting improvement is relative to a named benchmark, information set and evaluation population.

\textbf{Equilibrium and optimization.} An equilibrium comprises feasible plans, optimality, beliefs where needed, budget constraints and all market/resource clearing. The primitives or a stated benchmark define each correspondence \(\mathcal E(p;m)\). Nonemptiness is assumed or proved, never inferred just from compact policy sets. A selected-equilibrium target uses a declared selection rule; a robust target quantifies over all permitted equilibria. Use suprema or infima unless attainment is established. An \(\varepsilon\)-optimal policy is feasible and within \(\varepsilon\) of the finite optimum value. Report an empty feasible set separately.

\textbf{Welfare and accounting.} Normative utility, interpersonal normalization, social weights, numeraire, discounting and the use of public receipts are primitives. Behavioral utility need not coincide with the welfare criterion. Household welfare includes ownership income and rents once; adding profits again to compensating variation generally double-counts them. A monetary equivalent is defined by an explicit transfer and price convention, with monotonicity and range conditions. Average effects, mechanism contributions, transfers and welfare losses are different targets. Infinite sums require convergence and dynamic feasible plans require terminal or no-Ponzi conditions.

\textbf{Source versions.} A working-paper date, revision date and journal publication year are distinguished. A DOI for a journal version is not automatically the DOI of an earlier manuscript. Locators refer to the stated version and printed pages where specified. A metadata-only check does not verify a claim about a full-text conclusion. Every entry's source subsection narrows the provenance claim accordingly.
% END ECONOMICS STATEMENT
\end{document}
